\documentclass[aspectratio=169,xcolor=dvipsnames]{beamer}

\AtBeginSection[]
{
    \begin{frame}
        \frametitle{Visão Geral}
        \tableofcontents[currentsection]
    \end{frame}
}

\usefonttheme{professionalfonts} % using non standard fonts for beamer
\usefonttheme{serif} % default family is serif
\usepackage{fontspec}
\setmainfont{XCharter}[
    UprightFont    = *-Roman,
    BoldFont       = *-Bold,
    ItalicFont     = *-Italic,
    BoldItalicFont = *-BoldItalic,
]

\usepackage{hyperref}
\usepackage{graphicx} % Allows including images
\usepackage{booktabs} % Allows the use of \toprule, \midrule and \bottomrule in tables
\input{./slides/SimplePlusTheme/beamerthemeSimplePlus.sty}
\input{./slides/SimplePlusTheme/beamerinnerthemeSimplePlus.sty}
\input{./slides/SimplePlusTheme/beamerfontthemeSimplePlus.sty}
\input{./slides/SimplePlusTheme/beamercolorthemeSimplePlus.sty}

\usepackage{listings}
\setbeamercovered{transparent}

\usepackage{biblatex}
\addbibresource{slides/reference.bib}
\renewcommand*{\bibfont}{\footnotesize}
% Tamanho dos slides de "Leituras Recomendadas": maior que a bibliografia
% completa (\bibfont acima, que rende \tiny via a redefinição de \footnotesize)
% para destacar as leituras curadas.
\newcommand{\leiturasize}{\small}

\usepackage[brazilian ]{babel}
\usepackage{csquotes}
\usepackage{pgfplots}
\pgfplotsset{compat=1.18}

\usepackage{emoji}
\usepackage{amsmath}
\usepackage[xcharter,bigdelims,vvarbb]{newtxmath}
% newtxmath's operators font requests the fontspec family `XCharter(0)' in OT1;
% without these substitutions math digits and operator names silently fall back
% to Computer Modern (LaTeX Font Warning: Font shape `OT1/XCharter(0)/m/n' undefined).
\DeclareFontFamily{OT1}{XCharter(0)}{}
\DeclareFontShape{OT1}{XCharter(0)}{m}{n}{<->ssub * XCharter-TLF/m/n}{}
\DeclareFontShape{OT1}{XCharter(0)}{m}{it}{<->ssub * XCharter-TLF/m/it}{}
\DeclareFontShape{OT1}{XCharter(0)}{b}{n}{<->ssub * XCharter-TLF/b/n}{}
\DeclareFontShape{OT1}{XCharter(0)}{bx}{n}{<->ssub * XCharter-TLF/b/n}{}

\usepackage{bm}
\usepackage[table]{xcolor}
\usepackage{xcolor}
\usepackage{algpseudocode}

\usepackage{cancel}
\usepackage{ulem}

\usetikzlibrary{shapes}
\usetikzlibrary{arrows}
\usetikzlibrary {arrows.meta}
\usetikzlibrary{calc,positioning}
\usepgfplotslibrary{fillbetween}
\usetikzlibrary{angles,quotes}
\usetikzlibrary{tikzmark}


\definecolor{PastelRed}{HTML}{FFC1CC}
\definecolor{PastelBlue}{HTML}{C2F0FF}
\definecolor{PastelBlue2}{HTML}{3182BD}

\definecolor{PastelPink}{HTML}{FFCBE1}
\definecolor{PastelGreen}{HTML}{D6E5BD}
\definecolor{PastelYellow}{HTML}{F9E1A8}
\definecolor{PastelBlue}{HTML}{BCD8EC}
\definecolor{PastelPurple}{HTML}{DCCCEC}
\definecolor{PastelOrange}{HTML}{FFDAB4}


\renewcommand{\footnotesize}{\tiny}

\newcommand{\mespor}{
  \ifcase\month
    \or janeiro
    \or fevereiro
    \or março
    \or abril
    \or maio
    \or junho
    \or julho
    \or agosto
    \or setembro
    \or outubro
    \or novembro
    \or dezembro
  \fi
}
% \usepackage{csquotes}
% \usepackage{minted}
% \usemintedstyle{xcode}
\definecolor{vscLightBackground}{HTML}{FFFFFF}
\definecolor{vscLightText}{HTML}{000000}
\definecolor{vscLightGreen}{HTML}{008000}        % For Comments AND Numbers
\definecolor{vscLightRed}{HTML}{A31515}          % For Strings
\definecolor{vscLightTeal}{HTML}{267F99}
\definecolor{vscBlue}{HTML}{0000FF}% For main keywords (if, for, def...)
\definecolor{vscLightPytorch}{HTML}{800080}      % For PyTorch keywords (dark magenta)

\lstdefinestyle{vscode-light}{
    language=Python,
    backgroundcolor=\color{vscLightBackground},
    basicstyle=\ttfamily\scriptsize\color{vscLightText},
    keywordstyle=\color{vscLightTeal},
    commentstyle=\color{vscLightGreen},
    stringstyle=\color{vscLightRed},
    % Style for PyTorch keywords (using keyword group 2)
    keywordstyle=[2]\color{vscLightPytorch}, 
    morekeywords={self, True, False, None},
    morekeywords=[2]{
        torch, Tensor, device, dtype, to, nn, optim, autograd, utils, data,
        Module, Linear, Conv2d, ReLU, CrossEntropyLoss, Adam, SGD,
        Dataset, DataLoader, randn, zeros, ones, arange, cat, stack,
        view, reshape, sum, mean, max, min, matmul, no_grad, backward
    },
    keywordstyle=[3]\color{vscBlue}, 
    morekeywords=[3]{
        model, loss, y, y_hat, x , optimizer
    },
    frame=single,
    tabsize=1,
    showstringspaces=false,
    breaklines=true,
    captionpos=b,
    extendedchars=true,
    numbers=left,
}

\usepackage[cache=true]{minted}
% minted v3+ (TeX Live 2024+) auto-detects the output directory.
% minted v2 (Ubuntu apt TeX Live 2023) needs it set explicitly to
% match latexmkrc's $out_dir = 'build'.
\makeatletter
\@ifpackagelater{minted}{2024/01/01}{}{%
  \def\minted@outputdir{build/}%
}
\makeatother
\setminted{
    fontsize=\scriptsize,
    style=vs,
    breaklines=true,
    numbers=left,
    numbersep=5pt,
    frame=single,
}
\providecommand{\citeauthoryear}[1]{(\citeauthor{#1},~\citeyear{#1})}

\providecommand{\thetav}{\bm{\theta}}

% vetor coluna 2D com colchetes
\providecommand{\cvec}[2]{\begin{bmatrix} #1 \\ #2 \end{bmatrix}}

\providecommand{\varvector}[1]{\mathbf{#1}}
% instance
\providecommand{\X}{\varvector{X}}
\providecommand{\mlinstance}{\varvector{X}}
\providecommand{\mlinstancei}{\varvector{X}^{(i)}}
\providecommand{\mlinstancex}[1]{\varvector{X}^{(#1)}}

% target
\providecommand{\y}{\varvector{y}}
\providecommand{\mltarget}{\varvector{y}}
\providecommand{\mltargeti}{\varvector{y}^{(i)}}
\providecommand{\mltargetx}[1]{\varvector{y}^{(#1)}}

%hats
\providecommand{\yhat}{\hat{\varvector{y}}}
\providecommand{\mltargethat}{\hat{\varvector{y}}}
\providecommand{\mltargethati}{\hat{\varvector{y}}^{(i)}}
\providecommand{\mltargethatx}[1]{\hat{\varvector{y}}^{(#1)}}

% linear reg
\providecommand{\mllinearregression}{\mltargethat=\mlinstance\bm{\theta}}

\providecommand{\tikzarrowcolor}[3]{%
  \begin{tikzpicture}[remember picture, overlay]
    % Arrow body
    \path[fill=#3] ([shift={(#1,#2)}]current page.south west) 
      rectangle ++(0.3, -0.2);
    % Arrow head
    \path[fill=#3] ([shift={(#1+0.3,#2+0.1)}]current page.south west) -- 
         ++(0, -0.4) -- 
         ++(0.2, 0.2) -- 
         cycle;
  \end{tikzpicture}%
}

\providecommand{\tikzarrow}[2]{%
    \tikzarrowcolor{#1}{#2}{Red}
}



\DeclareMathOperator*{\argmax}{argmax}
\DeclareMathOperator*{\argmin}{argmin}


%------------------------------------------------------------------------
%   SHARED COLOR PALETTE FOR UNIT 05: GENERATIVE MODELS
%
%   Input this file (after preamble/custom-commands) in every lecture of
%   this unit so that the same concept always gets the same color.
%------------------------------------------------------------------------

% Encoder / Discriminator / "Going into the latent" path
\definecolor{EncRed}{HTML}{C62828}

% Decoder / Generator / "Coming out of the latent" path
\definecolor{DecBlue}{HTML}{1565C0}

% Aliases so either naming convention works (GANs call them Gen/Disc,
% VAEs/diffusion call them Encoder/Decoder).
\colorlet{GenRed}{EncRed}
\colorlet{DiscBlue}{DecBlue}

% Reconstruction / "true-data" anchor
\colorlet{ReconGreen}{MediumGreen}

% Noise / stochastic perturbation
\definecolor{NoiseAmber}{HTML}{F9A825}

% Accent for cross-references (prior lecture, forward pointers)
\definecolor{AccentBlue}{HTML}{1565C0}
\colorlet{FlowRed}{EncRed}
\colorlet{FlowGreen}{MediumGreen}

% Emphasis box color for pedagogical asides
\definecolor{SoftAmber}{HTML}{F9A825}

%------------------------------------------------------------------------
%   SHARED TIKZ: GENERATIVE-MODELING TRILEMMA
%
%   TikZ replacement for slides/img/05-geradores/generative-trilemma.png.
%   Input this file AFTER colors.tex in every lecture of unit 05.
%
%   Usage:  \resizebox{\linewidth}{!}{\generativetrilemma[<family>]}
%   where <family> in {gan, diffusion, vaenf, none} highlights the family
%   relevant to the current deck (default: none = neutral figure).
%------------------------------------------------------------------------

\usepackage{etoolbox} % \ifstrequal (no-op if beamer already loaded it)

\newcommand{\generativetrilemma}[1][none]{%
  % --- per-family stroke width and opacity, driven by the [#1] highlight ---
  \def\wGan{1pt}\def\oGan{1}%
  \def\wDif{1pt}\def\oDif{1}%
  \def\wVae{1pt}\def\oVae{1}%
  \ifstrequal{#1}{none}{}{%
    \def\oGan{0.30}\def\oDif{0.30}\def\oVae{0.30}%
    \ifstrequal{#1}{gan}{\def\oGan{1}\def\wGan{1.7pt}}{}%
    \ifstrequal{#1}{diffusion}{\def\oDif{1}\def\wDif{1.7pt}}{}%
    \ifstrequal{#1}{vaenf}{\def\oVae{1}\def\wVae{1.7pt}}{}%
  }%
  \begin{tikzpicture}[
      corner/.style={circle, draw=black, thick, minimum size=2.4cm,
                     align=center, font=\footnotesize},
      contour/.style={dashed, line cap=round},
    ]
    % --- three corner circles ---
    \node[corner] (T) at (0,2.4)   {High Quality\\Samples};
    \node[corner] (L) at (-2.5,-1.6) {Fast\\Sampling};
    \node[corner] (R) at (2.5,-1.6)  {Mode\\Coverage /\\Diversity};

    % --- family contours, each enclosing the two corners it reaches ---
    % GANs: quality + fast sampling (top + bottom-left)
    \draw[contour, DiscBlue, line width=\wGan, opacity=\oGan,
          rotate around={58:(-1.25,0.4)}]
          (-1.25,0.4) ellipse [x radius=3.9, y radius=1.5];
    % Diffusion: quality + coverage (top + bottom-right)
    \draw[contour, EncRed, line width=\wDif, opacity=\oDif,
          rotate around={-58:(1.25,0.4)}]
          (1.25,0.4) ellipse [x radius=3.9, y radius=1.5];
    % VAEs / Normalizing Flows: fast sampling + coverage (bottom edge)
    \draw[contour, NoiseAmber, line width=\wVae, opacity=\oVae]
          (0,-1.6) ellipse [x radius=4.0, y radius=1.5];

    % --- family labels, color-matched ---
    \node[DiscBlue, align=center, font=\small, opacity=\oGan]
      at (-4.9,2.2) {Generative\\Adversarial\\Networks};
    \node[EncRed, align=center, font=\small, opacity=\oDif]
      at (4.9,2.2) {Denoising\\Diffusion\\Models};
    \node[NoiseAmber, align=center, font=\small, opacity=\oVae]
      at (0,-3.6) {Variational Autoencoders,\\Normalizing Flows};
  \end{tikzpicture}%
}


%------------------------------------------------------------------------
%   TITLE PAGE
%------------------------------------------------------------------------

\title{Aprendizado Profundo}
\subtitle{\textit{Autoencoders} e \textit{Variational Autoencoders}}

\author{Lucas Silveira Kupssinskü}

\institute
{
    Machine Learning Theory and Applications Laboratory \\
    PUCRS
}
\date{\mespor de \the\year}

\begin{document}

  \begin{frame}
    \titlepage
  \end{frame}

  \begin{frame}{Visão Geral}
    \tableofcontents
  \end{frame}

  %======================================================================
  % SECTION 1: AUTOENCODERS
  %======================================================================
  \section{Autoencoders}

  %--- Slide: definição ---
  \begin{frame}{\textit{Autoencoders}}
    Um \textit{autoencoder} é uma rede neural treinada na tarefa de
    \textbf{reconstrução da própria entrada}.
    \vspace{0.4em}
    \begin{itemize}
      \item Uma das primeiras formas de aprendizado
      \textbf{auto-supervisionado} / não supervisionado.
      \item A ideia é criar uma \textbf{representação útil} das
      instâncias de interesse na etapa intermediária (\emph{gargalo}).
    \end{itemize}
    \begin{block}{Objetivo}
      \[
        \argmin_{\thetav,\bm{\phi}}
        J\!\left(\X,\, f\!\left(g(\X,\bm{\phi}),\thetav\right)\right)
      \]
      onde $g$ é o \textcolor{EncRed}{\textit{encoder}} (parâmetros $\bm{\phi}$)
      e $f$ é o \textcolor{DecBlue}{\textit{decoder}} (parâmetros $\thetav$).
    \end{block}
  \end{frame}

  %--- Slide: pipeline ---
  \begin{frame}{\textit{Autoencoders} -- Pipeline}
    \centering
    \includegraphics[width=0.82\linewidth]{slides/img/05-geradores/ae-encoder-decoder.png}\\
    \vspace{0.4em}
    \begin{itemize}
      \item \textcolor{EncRed}{\textit{Encoder}} $g$ reduz a entrada a
      uma representação compacta (o \textbf{latente}).
      \item \textcolor{DecBlue}{\textit{Decoder}} $f$ tenta reconstruir
      a entrada original a partir do latente.
    \end{itemize}
  \end{frame}

  %--- Slide: regularização ---
  \begin{frame}{\textit{Autoencoders} -- Regularização}
    \begin{itemize}
      \item É importante que o \textit{autoencoder} tenha
      \textbf{regularização} para evitar soluções triviais.
      \item Se \textbf{não houver um gargalo}, o \textit{autoencoder}
      pode simplesmente copiar a imagem.
      \item Se houver \textbf{unidades suficientes} no vetor latente,
      é possível \emph{indexar} as imagens originais (memorização).
    \end{itemize}
    \begin{alertblock}{Formas comuns de regularização}
      \begin{itemize}
        \item \textbf{Bottleneck}: dimensão latente $\ll$ dimensão da
        entrada.
        \item \textbf{Ruído} na entrada (denoising AE).
        \item \textbf{Esparsidade} na ativação do latente.
      \end{itemize}
    \end{alertblock}
  \end{frame}

  %--- Slide: denoising ---
  \begin{frame}{\textit{Denoising Autoencoders}\footfullcite{vincent2008denoising}}
    \begin{columns}[T]
      \begin{column}{0.48\textwidth}
        \begin{itemize}
          \item Alguns \textit{autoencoders} são treinados com o objetivo
          de \textbf{remover ruído} das instâncias de entrada.
          \item Tipicamente o \textit{noise} segue uma:
          \begin{itemize}
            \item distribuição \textbf{normal} (ruído aditivo);
            \item distribuição de \textbf{Bernoulli} (ruído
            multiplicativo que invalida algumas \textit{features}).
          \end{itemize}
          \item A reconstrução é avaliada contra a entrada
          \emph{original} (sem ruído).
        \end{itemize}
      \end{column}
      \begin{column}{0.50\textwidth}
        \centering
        \includegraphics[width=\linewidth]{slides/img/05-geradores/ae-denoising.png}
      \end{column}
    \end{columns}
  \end{frame}

  %--- Slide: aplicação colorização ---
  \begin{frame}{Aplicação: Colorização}
    \begin{columns}[T]
      \begin{column}{0.55\textwidth}
        \begin{itemize}
          \item \textbf{Entrada}: imagem P\&B (1 canal).
          \textbf{Saída}: a mesma imagem em cores (3 canais RGB).
          \item \textit{Encoder} comprime a cena em uma representação
          semântica; \textit{decoder} gera as cores plausíveis.
          \item Perda: MSE no espaço de cor (RGB, ou \textit{Lab}
          comparando apenas os canais $a, b$).
          \item Variação comum: \textit{U-Net} convolucional com
          \textit{skip connections} para preservar bordas.
        \end{itemize}
        \vspace{0.3em}
      \end{column}
      \begin{column}{0.43\textwidth}
        \centering
        \begin{tikzpicture}[scale=0.85]
          \node[draw, fill=black!10, minimum width=1.6cm, minimum height=1.6cm,
                font=\scriptsize] (in) at (0,0) {P\&B};
          \node[draw, EncRed, fill=EncRed!15, minimum width=0.7cm,
                minimum height=1.6cm, font=\tiny] (enc) at (1.7,0) {Enc.};
          \node[draw, fill=NoiseAmber!25, minimum width=0.55cm,
                minimum height=0.7cm, font=\tiny] (lat) at (2.7,0) {$\mathbf{z}$};
          \node[draw, DecBlue, fill=DecBlue!15, minimum width=0.7cm,
                minimum height=1.6cm, font=\tiny] (dec) at (3.7,0) {Dec.};
          \node[draw, fill=red!15, minimum width=1.6cm, minimum height=1.6cm,
                font=\scriptsize] (out) at (5.4,0) {Color};
          \draw[->,thick] (in) -- (enc);
          \draw[->,thick] (enc) -- (lat);
          \draw[->,thick] (lat) -- (dec);
          \draw[->,thick] (dec) -- (out);
        \end{tikzpicture}\\
        \vspace{0.4em}
        \scriptsize Esquema: 1 canal $\to$ representação $\to$ 3 canais.
      \end{column}
    \end{columns}
            \footnotesize Problema \textbf{mal-posto}: várias colorações
        são plausíveis. MSE básico tende a cores \emph{desaturadas}
        (média entre opções); variantes probabilísticas (VAE, GAN,
        difusão) geram cores mais vívidas.
  \end{frame}

  %--- Slide: aplicação anomaly detection ---
  \begin{frame}{Aplicação: Detecção de Anomalias}
    \begin{columns}[T]
      \begin{column}{0.50\textwidth}
        \begin{itemize}
          \item Instâncias \textbf{normais} reconstroem bem (erro baixo);
          \textbf{anomalias} reconstroem mal (erro alto).
          \item Score de anomalia:
          \[
            s(\mathbf{x}) = \lVert \mathbf{x} - f(g(\mathbf{x},\bm{\phi}),\thetav) \rVert^{2}
          \]
          \item Um \textit{threshold} sobre $s(\mathbf{x})$ separa as duas classes.
          \item Casos típicos: fraude em transações, defeitos industriais
          (\textit{visual inspection}), intrusão em redes, monitoramento
          de equipamentos.
        \end{itemize}
      \end{column}
      \begin{column}{0.48\textwidth}
        \centering
        \begin{tikzpicture}
          \begin{axis}[
            width=\linewidth, height=4.6cm,
            xlabel={\scriptsize erro de reconstrução $s(\mathbf{x})$},
            ylabel={\scriptsize densidade},
            xtick=\empty, ytick=\empty,
            ymin=0, ymax=1.2,
            domain=0:10, samples=120,
            axis lines=left, no markers,
            legend style={font=\tiny,at={(0.98,0.98)},anchor=north east,draw=none,fill=none},
          ]
            \addplot[thick,EncRed,fill=EncRed!25,fill opacity=0.65] {exp(-(x-2)^2/0.6)} \closedcycle;
            \addlegendentry{normal}
            \addplot[thick,NoiseAmber,fill=NoiseAmber!30,fill opacity=0.65] {0.55*exp(-(x-6.5)^2/2.5)} \closedcycle;
            \addlegendentry{anômalo}
            \draw[dashed,thick,black!65] (axis cs:4.2,0) -- (axis cs:4.2,1.1);
            \node[above,font=\tiny] at (axis cs:4.2,1.1) {\textit{threshold}};
          \end{axis}
        \end{tikzpicture}
        \vspace{0.2em}
        \begin{alertblock}{Por que funciona?}
          \scriptsize O AE só aprendeu a distribuição dos dados de
          treino, ele não tem como reconstruir o que nunca viu.
        \end{alertblock}
      \end{column}
    \end{columns}
  \end{frame}



  %======================================================================
  % SECTION 2: VARIATIONAL AUTOENCODERS
  %======================================================================
  \section{Variational Autoencoders}

  %--- Slide: ponte AE -> VAE ---
  \begin{frame}{Por que o \textit{autoencoder} comum não gera?}
    Para \textbf{gerar} bastaria amostrar um código $\mathbf{z}$ e decodificá-lo.
    No AE determinístico isso \textbf{falha}:
    \begin{columns}[T]
      \begin{column}{0.56\textwidth}
        \begin{itemize}
          \item O encoder envia cada $\mathbf{x}$ a um ponto $\mathbf{z}$, mas \textbf{não
          controla a distribuição} dos códigos.
          \item O latente fica com \textbf{buracos} (regiões sem dado) e
          aglomerados irregulares.
          \item Sem uma $p(\mathbf{z})$ conhecida de onde amostrar, um $\mathbf{z}$
          aleatório cai num buraco $\to$ decoder gera \textbf{lixo}.
        \end{itemize}
      \end{column}
      \begin{column}{0.42\textwidth}
        \centering
        \begin{tikzpicture}[scale=0.82,
          dot/.style={circle, inner sep=1pt}]
          % AE: scattered with gaps
          \draw[rounded corners, black!40] (0,0) rectangle (2.3,2.3);
          \node[font=\tiny] at (1.15,2.5) {AE: buracos};
          \foreach \x/\y in {0.3/1.9,0.5/2.0,0.4/1.7,1.95/0.4,2.0/0.65,
            1.75/0.5,0.4/0.4,0.65/0.5,1.9/1.95,2.0/1.75,0.35/1.1}{
            \node[dot, fill=EncRed] at (\x,\y) {};}
          % VAE: gaussian blob
          \begin{scope}[xshift=2.9cm]
            \draw[rounded corners, black!40] (0,0) rectangle (2.3,2.3);
            \node[font=\tiny] at (1.15,2.5) {VAE: $\mathcal{N}(0,\mathbf{I})$};
            \foreach \x/\y in {1.15/1.15,1.0/1.3,1.35/1.1,1.1/1.0,1.3/1.35,
              0.95/1.1,1.45/1.3,1.15/0.9,1.0/1.45,1.4/1.45,1.55/1.05,
              0.85/1.4,1.15/1.55,1.2/0.75}{
              \node[dot, fill=DecBlue] at (\x,\y) {};}
          \end{scope}
        \end{tikzpicture}
      \end{column}
    \end{columns}
    \begin{block}{A ideia do VAE}
      Forçar o encoder a mapear os dados a uma distribuição \textbf{contínua
      e conhecida} (gaussiana): então $\mathbf{z}\sim\mathcal{N}(0,\mathbf{I})$ decodificado
      gera dados novos plausíveis.
    \end{block}
  \end{frame}

  %--- Slide: definição VAE ---
  \begin{frame}{\textit{Variational Autoencoders}\footfullcite{kingma2014auto}}
    \begin{itemize}
      \item São modelos geradores \textbf{probabilísticos} cujo objetivo
      é aprender uma distribuição $p(\mathbf{x})$ sobre os dados.
      \item É um \textbf{modelo não linear de variável latente}.
      \begin{itemize}
        \item Depois do treinamento podemos usar apenas o decoder
        (amostragem) ou apenas o encoder (inferência).
      \end{itemize}
    \end{itemize}
    \begin{block}{Panorama}
      O VAE é a \textbf{fusão} de duas ideias:
      \begin{itemize}
        \item \textit{Autoencoder}: arquitetura encoder/decoder.
        \item \textbf{Inferência variacional}: aproximar a
        verossimilhança intratável via um limite inferior (ELBO).
      \end{itemize}
    \end{block}
    \vspace{0.2em}
    \footnotesize Próximo slide: arquitetura completa, para servir de
    \emph{mapa} ao longo da matemática.
  \end{frame}

  %--- Slide: arquitetura VAE (panorama, posicionado cedo) ---
  \begin{frame}{Arquitetura VAE -- Panorama}
    \begin{columns}[T]
      \begin{column}{0.40\textwidth}
        \begin{itemize}
          \item \textcolor{EncRed}{\textit{Encoder}} $g[\mathbf{x},\bm{\phi}]$ recebe
          $\mathbf{x}$ e produz $\boldsymbol{\mu}(\mathbf{x}),\, \boldsymbol{\Sigma}(\mathbf{x})$.
          \item Definimos
          $q(\mathbf{z}\mid \mathbf{x},\bm{\phi}) = \mathcal{N}(\boldsymbol{\mu}, \boldsymbol{\Sigma})$.
          \item Amostramos $\mathbf{z}^{*} \sim q(\mathbf{z}\mid \mathbf{x},\bm{\phi})$.
          \item \textcolor{DecBlue}{\textit{Decoder}} $f[\mathbf{z}^{*},\thetav]$
          gera os parâmetros de $p(\mathbf{x}\mid \mathbf{z}^{*},\thetav)$.
          \item Perda: $-\mathrm{ELBO} =$ reconstrução $+$ KL.
        \end{itemize}

      \end{column}
      \begin{column}{0.58\textwidth}
        \centering
        \resizebox{\linewidth}{!}{%
        \begin{tikzpicture}[
          every node/.style={font=\scriptsize},
          block/.style={draw, thick, minimum width=1.1cm, minimum height=0.9cm,
                        rounded corners=2pt},
          enc/.style={block, draw=EncRed, fill=EncRed!10, text=EncRed},
          dec/.style={block, draw=DecBlue, fill=DecBlue!10, text=DecBlue},
          arrow/.style={->, thick, >=stealth},
          ann/.style={font=\tiny, text=black!70, align=center},
          callout/.style={draw, dashed, rounded corners=1pt, fill=white,
                          inner sep=2pt, font=\tiny, align=center}
        ]
          % Main pipeline
          \node (x) {$\mathbf{x}$};
          \node[enc, right=0.35cm of x] (g) {$g[\mathbf{x},\bm{\phi}]$};
          \node[right=0.5cm of g] (q) {$q(\mathbf{z}\mid \mathbf{x},\bm{\phi})$};
          \node[draw, NoiseAmber, fill=NoiseAmber!15, minimum width=0.55cm,
                minimum height=0.55cm, right=0.5cm of q] (z) {$\mathbf{z}^{*}$};
          \node[dec, right=0.5cm of z] (f) {$f[\mathbf{z}^{*},\thetav]$};
          \node[right=0.4cm of f] (px) {$p(\mathbf{x}\mid \mathbf{z}^{*},\thetav)$};

          \draw[arrow] (x) -- (g);
          \draw[arrow] (g) -- (q) node[midway, above, ann] {$\boldsymbol{\mu}, \boldsymbol{\Sigma}$};
          \draw[arrow] (q) -- (z) node[midway, above, yshift=0.18cm, ann] {amostra};
          \draw[arrow] (z) -- (f);
          \draw[arrow] (f) -- (px);

          % Labels under boxes
          \node[ann, below=0.08cm of x] {entrada};
          \node[ann, below=0.08cm of g] {encoder};
          \node[ann, below=0.08cm of q] {posterior aprox.};
          \node[ann, below=0.08cm of z] {amostra};
          \node[ann, below=0.08cm of f] {decoder};
          \node[ann, below=0.08cm of px] {verossimilhança};

          % Loss callouts above
          \node[callout, draw=EncRed, text=EncRed, above=1.0cm of q]
            (kl) {variacional $\approx$ \textit{prior}\\$D_{KL}[q(\mathbf{z}\mid \mathbf{x},\bm{\phi})\,\|\,p(\mathbf{z})]$};
          \node[callout, draw=black!60, above=1.0cm of z, xshift=0.6cm]
            (loss) {função de perda\\$\mathrm{ELBO}[\thetav,\bm{\phi}]$};
          \node[callout, draw=DecBlue, text=DecBlue, above=1.0cm of px]
            (rec) {dado deve ter prob. alta\\$\log p(\mathbf{x}\mid \mathbf{z}^{*},\thetav)$};

          \draw[->, dashed, EncRed!70, thin] (kl) -- (q);
          \draw[->, dashed, DecBlue!70, thin] (rec) -- (px);
        \end{tikzpicture}}\\
        \vspace{0.3em}
        \tiny Adaptado de Prince (2023), Fig.~17.9.
      \end{column}
    \end{columns}
            \begin{block}{Use como mapa}
          \scriptsize Toda a matemática a seguir refere-se a esses
          mesmos blocos.
        \end{block}
  \end{frame}

  %--- Slide: modelos de variáveis latentes ---
  \begin{frame}{Modelos de variáveis latentes}
    \begin{itemize}
      \item Abordagem \textbf{indireta} para modelar uma distribuição
      $p(\mathbf{x})$ sobre uma variável multidimensional $\mathbf{x}$.
      \item Marginalização de uma probabilidade conjunta $p(\mathbf{x}, \mathbf{z})$:
    \end{itemize}
    \vspace{0.3em}
    \[
      p(\mathbf{x}) = \int p(\mathbf{x}, \mathbf{z})\, d\mathbf{z} = \int p(\mathbf{x} \mid \mathbf{z})\, p(\mathbf{z})\, d\mathbf{z}
    \]
    \vspace{0.3em}
    \begin{block}{Motivação}
      Distribuições relativamente simples $p(\mathbf{x} \mid \mathbf{z})$ e $p(\mathbf{z})$ podem,
      por marginalização, definir distribuições $p(\mathbf{x})$
      \textbf{complexas}.
    \end{block}
  \end{frame}

  %--- Slide: MoG como modelo latente ---
  \begin{frame}{Exemplo: Mistura de Gaussianas 1D}
    \begin{columns}[T]
      \begin{column}{0.48\textwidth}
        \begin{itemize}
          \item $z$ é uma variável latente \textbf{discreta}, com
          distribuição \emph{a priori} categórica:
          $p(z = n) = \lambda_n$.
          \item A verossimilhança segue uma normal:
          $p(x \mid z = n) = \mathcal{N}(\mu_n, \sigma_n^{2})$.
          \item A distribuição marginal:
          \[
            p(x) = \sum_n \lambda_n\, \mathcal{N}(\mu_n, \sigma_n^{2}).
          \]
        \end{itemize}
      \end{column}
      \begin{column}{0.50\textwidth}
        \centering
        \begin{tikzpicture}[font=\small,
          declare function={
            gauss(\x,\m,\s)=exp(-((\x-\m)^2)/(2*\s^2))/(\s*sqrt(2*pi));
          }]
          % --- painel superior: prior discreto p(z) = lambda_n ---
          \begin{axis}[
            name=pz,
            width=6.6cm, height=3.0cm,
            axis lines=left,
            ybar, bar width=13pt,
            xmin=0.3, xmax=3.7, ymin=0, ymax=0.55,
            xtick={1,2,3}, xticklabels={$1$,$2$,$3$},
            ytick=\empty,
            xlabel={$z$ (componente)}, ylabel={$p(z)=\lambda_n$},
            clip=false,
          ]
            \addplot[bar shift=0pt, fill=EncRed!65, draw=EncRed] coordinates {(1,0.45)};
            \addplot[bar shift=0pt, fill=NoiseAmber!75, draw=NoiseAmber!80!black] coordinates {(2,0.35)};
            \addplot[bar shift=0pt, fill=ReconGreen!60, draw=ReconGreen] coordinates {(3,0.20)};
          \end{axis}
          % --- painel inferior: marginal continua p(x) ---
          \begin{axis}[
            name=px,
            at={(pz.below south west)}, anchor=north west, yshift=-0.5cm,
            width=6.6cm, height=4.0cm,
            axis lines=left,
            xmin=-0.5, xmax=6.5, ymin=0, ymax=0.55,
            xlabel={$x$}, ylabel={$p(x)$},
            xtick=\empty, ytick=\empty,
            clip=false,
          ]
            % componentes ponderados lambda_n N(mu_n, sigma_n^2)
            \addplot[domain=-0.5:6.5, samples=120, EncRed, dashed, thick]
              {0.45*gauss(x,1.5,0.45)};
            \addplot[domain=-0.5:6.5, samples=120, NoiseAmber, dashed, thick]
              {0.35*gauss(x,3.2,0.7)};
            \addplot[domain=-0.5:6.5, samples=120, ReconGreen, dashed, thick]
              {0.20*gauss(x,4.6,0.5)};
            % marginal = soma ponderada das componentes
            \addplot[domain=-0.5:6.5, samples=160, DecBlue, very thick,
                     fill=DecBlue, fill opacity=0.12]
              {0.45*gauss(x,1.5,0.45)+0.35*gauss(x,3.2,0.7)+0.20*gauss(x,4.6,0.5)}
              \closedcycle;
          \end{axis}
        \end{tikzpicture}\\
        \vspace{0.1em}
        \scriptsize O \textit{prior} discreto $p(z)$ pesa cada componente,
        e a marginal \textcolor{DecBlue}{$p(x)$} é a soma ponderada das
        gaussianas (tracejadas).
      \end{column}
    \end{columns}
  \end{frame}

  %--- Slide: modelos não lineares ---
  \begin{frame}{Modelos não lineares de variáveis latentes}
    \begin{columns}[T]
      \begin{column}{0.58\textwidth}
        Tanto $\mathbf{z}$ quanto $\mathbf{x}$ são \textbf{contínuos} e
        \textbf{multivariados}.
        \begin{itemize}
          \item \textit{Prior} normal padrão:
          $p(\mathbf{z}) = \mathcal{N}_{\mathbf{z}}(0, \mathbf{I})$.
          \item Verossimilhança gaussiana com média não linear
          $f(\mathbf{z}, \thetav)$ e covariância esférica:
          \[
            p(\mathbf{x} \mid \mathbf{z}, \thetav) = \mathcal{N}_{\mathbf{x}}\!\left(f(\mathbf{z}, \thetav), \sigma^{2}\mathbf{I}\right).
          \]
          \item Marginal:
          \[
            p(\mathbf{x} \mid \thetav) = \int p(\mathbf{x}\mid \mathbf{z},\thetav)\, p(\mathbf{z})\, d\mathbf{z}.
          \]
          \item Integrando gaussiano, mas $f$ não linear $\Rightarrow$
          integral \textbf{sem forma fechada}.
        \end{itemize}
      \end{column}
      \begin{column}{0.40\textwidth}
        \centering
        \begin{tikzpicture}[scale=0.55]
          \fill[EncRed!25,draw=EncRed,thick] (-2.6,0) circle (1.2);
          \node[font=\scriptsize] at (-2.6,-1.7) {$p(\mathbf{z})=\mathcal{N}(0,\mathbf{I})$};
          \node[font=\tiny,EncRed] at (-2.6,0) {simples};
          \draw[->,very thick,black!70] (-1.2,0) -- (1.0,0);
          \node[font=\scriptsize,above] at (-0.1,0.1) {$f(\mathbf{z},\thetav)$};
          \begin{scope}[shift={(2.6,0)}]
            \fill[DecBlue!25,draw=DecBlue,thick,smooth cycle]
              plot coordinates {(-1.1,0.5) (-0.5,1.2) (0.3,1.3) (1.0,0.8) (1.2,0.0) (0.7,-0.9) (-0.2,-1.1) (-0.9,-0.7) (-1.2,-0.1)};
            \node[font=\scriptsize] at (0,-1.7) {$p(\mathbf{x}\mid\thetav)$};
            \node[font=\tiny,DecBlue] at (0,0) {complexa};
          \end{scope}
        \end{tikzpicture}\\
        \vspace{0.6em}
        \scriptsize Função $f$ não linear transforma uma gaussiana
        simples em uma distribuição arbitrariamente complexa.
      \end{column}
    \end{columns}
  \end{frame}

  %--- Slide: ancestral sampling ---
  \begin{frame}{Gerando novas instâncias $\mathbf{x}^{*}$}
    Usamos \textbf{amostragem ancestral}:
    \vspace{0.3em}
    \begin{columns}[T]
      \begin{column}{0.46\textwidth}
        \begin{enumerate}
          \item Amostrar um latente $\mathbf{z}^{*}$ de $p(\mathbf{z})$
          (\textit{prior}).
          \item Passar $\mathbf{z}^{*}$ pela rede
          $f(\mathbf{z}^{*}, \thetav)$ para obter a média da
          verossimilhança $p(\mathbf{x} \mid \mathbf{z}^{*}, \thetav)$.
          \item Amostrar $\mathbf{x}^{*}$ de
          $p(\mathbf{x} \mid \mathbf{z}^{*}, \thetav)$.
        \end{enumerate}
        \vspace{0.4em}
        \footnotesize A rede $f(\mathbf{z},\thetav)$ é exatamente o
        \textcolor{DecBlue}{\textit{decoder}} da arquitetura.
      \end{column}
      \begin{column}{0.52\textwidth}
        \centering
        \resizebox{\linewidth}{!}{%
        \begin{tikzpicture}[font=\small,>=latex]
          % --- 1) espaço latente: prior p(z) ---
          \node[font=\scriptsize,black!60] at (0,1.5) {espaço latente $\mathbf{z}$};
          \fill[black!7] (0,0) circle (0.95);
          \draw[black!40,dashed] (0,0) circle (0.95);
          \foreach \p in {(-0.4,0.3),(0.2,-0.4),(-0.2,-0.15),(0.5,0.5),(-0.5,-0.35)}
            \fill[black!35] \p circle (0.9pt);
          \fill[EncRed] (0.35,0.28) circle (2pt);
          \node[EncRed,font=\scriptsize,above right=-1pt] at (0.35,0.28) {$\mathbf{z}^*$};
          \node[font=\scriptsize,black!75] at (0,-1.45)
            {\textbf{1.}~$\mathbf{z}^*\!\sim p(\mathbf{z})$};
          % --- 2) decoder f ---
          \draw[->,thick,black!70] (1.0,0) -- (1.72,0);
          \fill[DecBlue!16,draw=DecBlue,thick]
            (1.78,0.55) -- (3.08,1.0) -- (3.08,-1.0) -- (1.78,-0.55) -- cycle;
          \node[DecBlue,font=\scriptsize] at (2.46,0.12) {$f(\mathbf{z},\thetav)$};
          \node[DecBlue,font=\tiny] at (2.46,-0.3) {\textit{decoder}};
          \node[font=\scriptsize,black!75] at (2.43,-1.45) {\textbf{2.}};
          \draw[->,thick,black!70] (3.08,0) -- (3.8,0);
          % --- 3) espaço de dados: likelihood p(x|z*) ---
          \node[font=\scriptsize,black!60] at (4.65,1.5) {espaço de dados $\mathbf{x}$};
          \fill[DecBlue!10] (4.65,0) ellipse (0.95 and 0.78);
          \draw[DecBlue!45] (4.65,0) ellipse (0.95 and 0.78);
          \draw[DecBlue!60] (4.65,0) ellipse (0.55 and 0.45);
          \node[DecBlue,font=\scriptsize] at (4.65,0) {$\boldsymbol{\times}$};
          \node[font=\tiny,DecBlue,below=0pt] at (4.65,-0.06) {$\boldsymbol{\mu}\!=\!f(\mathbf{z}^*,\thetav)$};
          \fill[NoiseAmber!90!black] (5.05,0.34) circle (2pt);
          \node[font=\scriptsize,above right=-1pt] at (5.05,0.34) {$\mathbf{x}^*$};
          \node[font=\scriptsize,black!75] at (4.65,-1.45)
            {\textbf{3.}~$\mathbf{x}^*\!\sim p(\mathbf{x}\mid \mathbf{z}^*,\thetav)$};
        \end{tikzpicture}%
        }\\
        \vspace{0.2em}
        \scriptsize Amostramos $\mathbf{z}^*$ do \textit{prior}, decodificamos com
        $f(\mathbf{z},\thetav)$ e amostramos $\mathbf{x}^*$ da verossimilhança.\\
        \textcolor{black!55}{Na prática, costuma-se usar a média
        $f(\mathbf{z}^*,\thetav)$ diretamente, sem somar o ruído $\sigma^2\mathbf{I}$.}
      \end{column}
    \end{columns}
  \end{frame}

  %======================================================================
  % SECTION 3: TREINAMENTO E ELBO
  %======================================================================
  \section{Treinamento e ELBO}

  %--- Slide: Treinamento intratável ---
  \begin{frame}{Treinamento -- O problema}
    Maximizar a log-verossimilhança direta é \textbf{inviável}:
    \[
      \hat{\thetav} = \argmax_{\thetav}
      \left[\sum_{i=1}^{N} \int
        \mathcal{N}_{\mathbf{x}^{(i)}}\!\left(f(\mathbf{z}, \thetav), \sigma^{2}\mathbf{I}\right)
        \mathcal{N}_{\mathbf{z}}(0, \mathbf{I})\, d\mathbf{z}
      \right]
    \]
    \vspace{0.2em}
    \begin{itemize}
      \item Integral \textbf{sem forma fechada}, pois $f$ é não linear.
      \item Monte Carlo direto $p(\mathbf{x}) \approx \tfrac{1}{K}\sum_{k} p(\mathbf{x}\mid \mathbf{z}_k)$
      com $\mathbf{z}_k \sim p(\mathbf{z})$ também \textbf{falha}: para um $\mathbf{x}$ específico,
      $\mathbf{z}$ amostrado do prior quase nunca explica $\mathbf{x}$, então
      $p(\mathbf{x}\mid \mathbf{z})$ é minúsculo na enorme maioria das amostras.
    \end{itemize}
    \begin{alertblock}{Estratégia}
      Maximizamos um \textbf{limite inferior} da log-verossimilhança,
      usando uma proposta inteligente para amostrar $\mathbf{z}$.
    \end{alertblock}
  \end{frame}

  %--- Slide: motivação do q via importance sampling ---
  \begin{frame}{Por que precisamos de um \textit{encoder}?}
    \textbf{Ideia}: ao invés de amostrar $\mathbf{z}$ do prior, sortear $\mathbf{z}$ de
    uma distribuição que \emph{concentra massa onde $\mathbf{z}$ explica $\mathbf{x}$}.
    \vspace{0.4em}
    \begin{itemize}
      \item Essa distribuição ideal é a posterior $p(\mathbf{z}\mid \mathbf{x})$, mas ela
      também é intratável.
      \item \textbf{Solução}: aproximamos com uma proposta parametrizada
      \[
        q(\mathbf{z}\mid \mathbf{x}, \bm{\phi}) \approx p(\mathbf{z}\mid \mathbf{x}),
      \]
      implementada por uma \textbf{segunda rede} com parâmetros $\bm{\phi}$.
      \item Essa rede é o \textcolor{EncRed}{\textit{encoder}} do VAE:
      para cada $\mathbf{x}$, produz uma distribuição sobre $\mathbf{z}$.
    \end{itemize}
    \vspace{0.2em}
    \begin{block}{O que ganhamos}
      Com $q$ no lugar de $p(\mathbf{z})$, podemos derivar um \textbf{limite
      inferior} (ELBO) que é \emph{apertado} quando $q$ aproxima bem a
      posterior, vamos ver isso a seguir.
    \end{block}
  \end{frame}

  %--- Slide: prior amplo vs proposta concentrada ---
  \begin{frame}{Amostrar do \textit{prior} é ineficiente}
    \begin{columns}[c]
      \begin{column}{0.58\textwidth}
        \centering
        \begin{tikzpicture}[scale=1.0,>=latex]
          % eixos latentes
          \draw[->,black!45] (-2.9,0) -- (2.9,0) node[right,font=\tiny,black!55] {$z_1$};
          \draw[->,black!45] (0,-2.6) -- (0,2.6) node[above,font=\tiny,black!55] {$z_2$};
          % prior p(z) = N(0,I): massa espalhada
          \fill[black!7] (0,0) circle (2.4);
          \draw[black!45,dashed,thick] (0,0) circle (2.4);
          \draw[black!30,dashed] (0,0) circle (1.5);
          \node[font=\scriptsize,black!60] at (-1.55,-2.05) {$p(\mathbf{z})=\mathcal{N}(0,\mathbf{I})$};
          % amostras do prior (massa espalhada)
          \foreach \p in {%
            (-2.0,0.3),(-1.6,1.1),(-1.7,-0.8),(-1.2,1.7),(-1.3,-1.5),%
            (-0.9,0.6),(-0.7,-0.4),(-0.5,1.3),(-0.4,-1.6),(-0.2,0.2),%
            (0.0,1.0),(0.1,-0.9),(0.3,1.8),(0.5,-1.4),(0.6,0.4),%
            (0.8,-0.5),(0.9,1.5),(1.1,-1.1),(1.3,0.2),(1.5,-0.7),%
            (1.8,-1.3),(2.0,0.6),(2.1,-0.2),(-2.2,-0.6),(-0.1,-1.9),%
            (1.0,-1.8),(2.0,1.6),(-1.9,1.4)}
            \fill[black!45] \p circle (1.6pt);
          % proposta concentrada q(z|x,phi)
          \fill[EncRed!22] (1.6,1.1) ellipse (0.5 and 0.42);
          \draw[EncRed,thick] (1.6,1.1) ellipse (0.5 and 0.42);
          % amostras que de fato explicam x (dentro de q)
          \foreach \p in {(1.5,1.02),(1.72,1.18),(1.6,0.95)}
            \fill[ReconGreen] \p circle (2.0pt);
          \node[font=\scriptsize,EncRed,align=center] at (1.6,2.05)
            {$q(\mathbf{z}\mid \mathbf{x},\bm{\phi})$};
        \end{tikzpicture}
      \end{column}
      \begin{column}{0.40\textwidth}
        \footnotesize
        \begin{itemize}
          \item \textcolor{black!60}{\textbf{Cinza}}: amostras do
          \textit{prior} $p(\mathbf{z})$. A massa fica \emph{espalhada} por todo
          o espaço latente.
          \item \textcolor{ReconGreen}{\textbf{Verde}}: as raras amostras
          que caem onde $\mathbf{z}$ realmente \emph{explica} $\mathbf{x}$, isto é, onde
          $p(\mathbf{x}\mid \mathbf{z})$ é alto.
          \item Acertar essa região por sorteio do \textit{prior} exige
          um número \textbf{enorme} de amostras: $p(\mathbf{z})$ é pouco
          informativo sobre um $\mathbf{x}$ específico.
        \end{itemize}
        \vspace{0.3em}
        \begin{block}{Solução}
          A proposta \textcolor{EncRed}{$q(\mathbf{z}\mid \mathbf{x},\bm{\phi})$} já
          \textbf{concentra} a massa nessa região.
        \end{block}
      \end{column}
    \end{columns}
  \end{frame}

  %--- Slide: Jensen ---
  \begin{frame}{Desigualdade de Jensen}
    \begin{columns}[c]
      \begin{column}{0.38\textwidth}
        Para qualquer função \textbf{côncava} $g(\cdot)$:
        \[
          g\!\left(\mathbb{E}[x]\right) \geq \mathbb{E}\!\left[g(x)\right]
        \]
        Caso da função log (com $\log Y \sim \mathcal{N}(\mu,\sigma^2)$):
        \[
          \log\!\left(\mathbb{E}[Y]\right) \geq \mathbb{E}\!\left[\log Y\right]
        \]
        \footnotesize Aqui $\mu=0$, $\sigma=1$, então a folga
        $\sigma^2/2 = 0.5$.
      \end{column}
      \begin{column}{0.60\textwidth}
        \centering
        \resizebox{\linewidth}{!}{%
        \begin{tikzpicture}[font=\small,
          declare function={
            gauss(\x,\m,\s)=exp(-((\x-\m)^2)/(2*\s^2))/(\s*sqrt(2*pi));
            lognorm(\y,\m,\s)=exp(-((ln(\y)-\m)^2)/(2*\s^2))/(\y*\s*sqrt(2*pi));
          }]
          % ---- ESQUERDA: Normal de log Y ----
          \begin{axis}[
            name=left,
            width=5.4cm, height=4.3cm,
            axis lines=left,
            xmin=-3, xmax=3, ymin=0, ymax=0.58,
            xlabel={$\log y$}, ylabel={densidade},
            xtick=\empty, ytick=\empty,
            title={\scriptsize $\log Y \sim$ Normal},
            clip=false,
          ]
            \addplot[domain=-3:3, samples=140, black!55, thick,
                     fill=black!8] {gauss(x,0,1)} \closedcycle;
            \draw[ReconGreen, very thick] (axis cs:0,0) -- (axis cs:0,0.52);
            \node[ReconGreen, font=\tiny, above] at (axis cs:0,0.52)
              {$\mathbb{E}[\log Y]=0$};
            \draw[EncRed, very thick, dashed] (axis cs:0.5,0) -- (axis cs:0.5,0.40);
            \node[EncRed, font=\tiny, anchor=west] at (axis cs:0.62,0.36)
              {$\log\mathbb{E}[Y]=0.5$};
          \end{axis}
          % ---- DIREITA: Log-normal de Y ----
          \begin{axis}[
            name=right,
            at={(left.east)}, xshift=1.1cm, anchor=west,
            width=5.4cm, height=4.3cm,
            axis lines=left,
            xmin=0, xmax=5, ymin=0, ymax=0.78,
            xlabel={$y$}, ylabel={densidade},
            xtick=\empty, ytick=\empty,
            title={\scriptsize $Y \sim$ Log-normal},
            clip=false,
          ]
            \addplot[domain=0.02:5, samples=200, black!55, thick,
                     fill=black!8] {lognorm(x,0,1)} \closedcycle;
            \draw[ReconGreen, very thick] (axis cs:1,0) -- (axis cs:1,0.70);
            \node[ReconGreen, font=\tiny, above] at (axis cs:1,0.70)
              {$e^{\mathbb{E}[\log Y]}=1$};
            \draw[EncRed, very thick, dashed] (axis cs:1.65,0) -- (axis cs:1.65,0.56);
            \node[EncRed, font=\tiny, anchor=west] at (axis cs:1.78,0.52)
              {$\mathbb{E}[Y]\approx 1.65$};
          \end{axis}
        \end{tikzpicture}%
        }\\
        \vspace{0.2em}
        \scriptsize A cauda da log-normal puxa a média aritmética
        \textcolor{EncRed}{$\mathbb{E}[Y]$} à direita da geométrica
        \textcolor{ReconGreen}{$e^{\mathbb{E}[\log Y]}$}; logo
        $\log\mathbb{E}[Y] \geq \mathbb{E}[\log Y]$.
      \end{column}
    \end{columns}
  \end{frame}

  %--- Slide: refresher ELBO ---
  \begin{frame}{De onde vem o ELBO?}
    O nome \textbf{ELBO} (\textit{Evidence Lower BOund}) é herdado de
    \textbf{inferência variacional}.
    \vspace{0.3em}
    \begin{itemize}
      \item \textbf{Evidência}: na linguagem bayesiana, $p(\mathbf{x}\mid\thetav)
      = \int p(\mathbf{x},\mathbf{z}\mid\thetav)\,d\mathbf{z}$ é a \emph{verossimilhança marginal},
      a ``evidência'' que os dados dão sobre o modelo.
      \item \textbf{Limite inferior}: como a integral é intratável,
      construímos uma função $\mathcal{L}[\thetav,\bm{\phi}]$ tal que
      $\log p(\mathbf{x}\mid\thetav) \geq \mathcal{L}[\thetav,\bm{\phi}]$, e
      maximizamos $\mathcal{L}$ no lugar de $\log p(\mathbf{x})$.
    \end{itemize}
    \begin{block}{Ingrediente que falta}
      \textbf{Importance sampling}: para qualquer $q$ com suporte sobre
      $\mathbf{z}$,
      $p(\mathbf{x}\mid\thetav) = \mathbb{E}_{\mathbf{z}\sim q}\!\left[p(\mathbf{x},\mathbf{z}\mid\thetav)/q(\mathbf{z})\right]$.
      O próximo slide combina esse passo com Jensen para obter o limite
      inferior.
    \end{block}
  \end{frame}

  %--- Slide: derivação do limite inferior ---
  \begin{frame}{Limite Inferior -- Derivação}
    \small
    Reescrevemos a integral como uma esperança sob $q$ e aplicamos
    Jensen na concavidade do $\log$:
    {\small
    \begin{align*}
      \log p(\mathbf{x} \mid \thetav)
      &= \log\!\left[\int p(\mathbf{x}, \mathbf{z} \mid \thetav)\, d\mathbf{z}\right]
         && \text{\scriptsize marginal} \\
      &= \log\!\left[\int q(\mathbf{z})\,\frac{p(\mathbf{x}, \mathbf{z} \mid \thetav)}{q(\mathbf{z})}\, d\mathbf{z}\right]
         && \text{\scriptsize multiplica e divide por $q(\mathbf{z})$} \\
      &= \log\!\left(\mathbb{E}_{\mathbf{z}\sim q}\!\left[\frac{p(\mathbf{x}, \mathbf{z} \mid \thetav)}{q(\mathbf{z})}\right]\right)
         && \text{\scriptsize integral ponderada por $q$ $=$ esperança} \\
      &\geq \mathbb{E}_{\mathbf{z}\sim q}\!\left[\log\frac{p(\mathbf{x}, \mathbf{z} \mid \thetav)}{q(\mathbf{z})}\right]
         && \text{\scriptsize Jensen ($\log$ côncavo)} \\
      &= \int q(\mathbf{z})\,\log\!\left(\frac{p(\mathbf{x}, \mathbf{z} \mid \thetav)}{q(\mathbf{z})}\right) d\mathbf{z}
         && \text{\scriptsize de volta à forma integral}
    \end{align*}}
    \begin{block}{ELBO}
      \footnotesize O $q(\mathbf{z})$ \emph{externo} é a medida da esperança,
      permanece como peso; apenas o $q(\mathbf{z})$ do denominador entra no
      $\log$. Esse limite é o \textit{evidence lower bound}; na prática
      $q$ é parametrizado: $q(\mathbf{z} \mid \mathbf{x},\bm{\phi})$.
    \end{block}
  \end{frame}

  %--- Slide: ponte ELBO -> identidade ---
  \begin{frame}{Limite Inferior -- Conexão com a posterior}
    \small
    A partir daqui escrevemos $q$ como $q(\mathbf{z}\mid \mathbf{x},\bm{\phi})$ (o encoder
    condiciona em $\mathbf{x}$). Aplicando a regra do produto
    $p(\mathbf{x},\mathbf{z}\mid\thetav) = p(\mathbf{z}\mid \mathbf{x},\thetav)\,p(\mathbf{x}\mid\thetav)$ no ELBO:
    {\small
    \begin{align*}
      \mathrm{ELBO}[\thetav,\bm{\phi}]
      &= \mathbb{E}_{q}\!\left[\log\tfrac{p(\mathbf{x},\mathbf{z}\mid\thetav)}{q(\mathbf{z}\mid \mathbf{x},\bm{\phi})}\right]
         && \text{\scriptsize definição} \\
      &= \mathbb{E}_{q}\!\left[\log\tfrac{p(\mathbf{z}\mid \mathbf{x},\thetav)\,p(\mathbf{x}\mid\thetav)}{q(\mathbf{z}\mid \mathbf{x},\bm{\phi})}\right]
         && \text{\scriptsize regra do produto} \\
      &= \mathbb{E}_{q}\!\left[\log p(\mathbf{x}\mid\thetav)\right]
       + \mathbb{E}_{q}\!\left[\log\tfrac{p(\mathbf{z}\mid \mathbf{x},\thetav)}{q(\mathbf{z}\mid \mathbf{x},\bm{\phi})}\right]
         && \text{\scriptsize separa o $\log$} \\
      &= \log p(\mathbf{x}\mid\thetav)
       - \mathbb{E}_{q}\!\left[\log\tfrac{q(\mathbf{z}\mid \mathbf{x},\bm{\phi})}{p(\mathbf{z}\mid \mathbf{x},\thetav)}\right]
         && \text{\scriptsize $\log p(\mathbf{x})$ sai da $\mathbb{E}$; troca o sinal} \\
      &= \log p(\mathbf{x}\mid\thetav)
       - D_{KL}\!\left[q(\mathbf{z}\mid \mathbf{x},\bm{\phi})\,\|\,p(\mathbf{z}\mid \mathbf{x},\thetav)\right]
         && \text{\scriptsize definição de KL}
    \end{align*}}
    \begin{alertblock}{Rearranjando}
      \footnotesize
      $\log p(\mathbf{x}\mid\thetav) = \mathrm{ELBO}[\thetav,\bm{\phi}]
       + D_{KL}\!\left[q(\mathbf{z}\mid \mathbf{x},\bm{\phi})\,\|\,p(\mathbf{z}\mid \mathbf{x},\thetav)\right]$,
      a identidade exata do próximo slide.
    \end{alertblock}
  \end{frame}

  %--- Slide: identidade ELBO + KL ---
  \begin{frame}{ELBO -- A identidade mestre}
    Da derivação anterior, chega-se à \textbf{identidade exata}:
    \vspace{0.3em}
    \[
      \boxed{\;\log p(\mathbf{x}\mid\thetav)
        = \underbrace{\int q(\mathbf{z}\mid \mathbf{x},\bm{\phi})\,\log \tfrac{p(\mathbf{x},\mathbf{z}\mid\thetav)}{q(\mathbf{z}\mid \mathbf{x},\bm{\phi})}\, d\mathbf{z}}_{\mathrm{ELBO}[\thetav,\bm{\phi}]}
        \;+\; \underbrace{D_{KL}\!\left[q(\mathbf{z}\mid \mathbf{x},\bm{\phi})\,\|\,p(\mathbf{z}\mid \mathbf{x},\thetav)\right]}_{\geq 0}\;}
    \]
    \vspace{0.2em}
    \begin{itemize}
      \item O \emph{gap} entre $\log p(\mathbf{x})$ e o ELBO é \textbf{exatamente}
      a KL entre $q$ e a posterior verdadeira.
      \item Bound \textbf{apertado} $\Leftrightarrow$ $q(\mathbf{z}\mid \mathbf{x},\bm{\phi})
      = p(\mathbf{z}\mid \mathbf{x},\thetav)$.
      \item Maximizar o ELBO faz \textbf{duas coisas ao mesmo tempo}:
      sobe $\log p(\mathbf{x})$ (aprende o modelo) \emph{e} aperta $q$ contra a
      posterior (aprende o encoder).
    \end{itemize}
  \end{frame}

  %--- Slide: ELBO reorganizado ---
  \begin{frame}{ELBO -- Reorganizando}
    \[
      \mathrm{ELBO}[\thetav, \bm{\phi}] =
      \int q(\mathbf{z} \mid \mathbf{x},\bm{\phi}) \log\!\left(\frac{p(\mathbf{x}, \mathbf{z} \mid \thetav)}{q(\mathbf{z} \mid \mathbf{x},\bm{\phi})}\right) d\mathbf{z}
    \]
    \vspace{0.2em}
    Aplicando $p(\mathbf{x}, \mathbf{z}) = p(\mathbf{x} \mid \mathbf{z})\,p(\mathbf{z})$ e separando a razão:
    \begin{align*}
      &= \int q(\mathbf{z} \mid \mathbf{x},\bm{\phi}) \log\!\left(\frac{p(\mathbf{x} \mid \mathbf{z}, \thetav)\,p(\mathbf{z})}{q(\mathbf{z} \mid \mathbf{x},\bm{\phi})}\right) d\mathbf{z} \\
      &= \int q(\mathbf{z} \mid \mathbf{x},\bm{\phi}) \log p(\mathbf{x} \mid \mathbf{z}, \thetav)\, d\mathbf{z}
      - D_{KL}\!\left[q(\mathbf{z} \mid \mathbf{x},\bm{\phi})\,\|\,p(\mathbf{z})\right]
    \end{align*}
  \end{frame}

  %--- Slide: interpretação ELBO ---
  \begin{frame}{ELBO -- Interpretação}
    \[
      \mathrm{ELBO}[\thetav, \bm{\phi}] =
      \underbrace{\int q(\mathbf{z} \mid \mathbf{x},\bm{\phi}) \log p(\mathbf{x} \mid \mathbf{z}, \thetav)\, d\mathbf{z}}_{\textcolor{DecBlue}{\text{Reconstruction Loss}}}
      \underbrace{- D_{KL}\!\left[q(\mathbf{z} \mid \mathbf{x},\bm{\phi})\,\|\,p(\mathbf{z})\right]}_{\textcolor{EncRed}{\text{Regularização do latente}}}
    \]
    \vspace{0.3em}
    \begin{itemize}
      \item \textcolor{DecBlue}{\textbf{Primeiro termo}}:
      log-verossimilhança esperada da reconstrução, sob $\mathbf{z}$ amostrados
      do encoder. Também chamado de \textit{reconstruction loss}.
      \item \textcolor{EncRed}{\textbf{Segundo termo}}: avalia quão
      perto a distribuição auxiliar $q(\mathbf{z} \mid \mathbf{x},\bm{\phi})$ está do
      \textit{prior} $p(\mathbf{z})$.
    \end{itemize}
    \vspace{0.2em}
    \begin{alertblock}{Sinal}
      O ELBO é \textbf{maximizado} (é um limite \emph{inferior} da
      log-verossimilhança); a função de custo treinada é
      $-\mathrm{ELBO}$, minimizada.
    \end{alertblock}
  \end{frame}

  %--- Slide: Reconstruction loss = MSE ---
  \begin{frame}{Termo de reconstrução $=$ MSE (caso gaussiano)}
    Para $p(\mathbf{x}\mid \mathbf{z},\thetav) = \mathcal{N}(f(\mathbf{z},\thetav),\,
    \sigma^{2}\mathbf{I})$:
    \vspace{0.2em}
    \begin{align*}
      \log p(\mathbf{x}\mid \mathbf{z},\thetav)
      &= \textcolor{ForestGreen}{-\tfrac{D_x}{2}\log(2\pi\sigma^{2})}
       -\tfrac{1}{2\sigma^{2}}\,\lVert \mathbf{x} - f(\mathbf{z},\thetav)\rVert^{2} \\[0.2em]
      &= -\tfrac{1}{2\sigma^{2}}\,\lVert \mathbf{x} - f(\mathbf{z},\thetav)\rVert^{2}
       + \textcolor{ForestGreen}{\mathrm{const}}
    \end{align*}
    \vspace{0.1em}
    \begin{block}{Ponte com o AE clássico}
      Maximizar $\log p(\mathbf{x}\mid \mathbf{z},\thetav)$ \textbf{é exatamente}
      minimizar o erro quadrático $\lVert \mathbf{x} - f(\mathbf{z},\thetav)\rVert^{2}$.
      O ELBO se decompõe em
      \[
        -\mathrm{ELBO} \;=\;
        \underbrace{\tfrac{1}{2\sigma^{2}}\mathbb{E}_{q}\!\left[\lVert \mathbf{x} - f(\mathbf{z},\thetav)\rVert^{2}\right]}_{\textcolor{DecBlue}{\text{reconstrução (MSE)}}}
        + \underbrace{D_{KL}\!\left[q(\mathbf{z}\mid \mathbf{x},\bm{\phi})\,\|\,p(\mathbf{z})\right]}_{\textcolor{EncRed}{\text{regularização}}}
        + \mathrm{const}.
      \]
    \end{block}
  \end{frame}

  %--- Slide: KL entre gaussianas 1/2 ---
  \begin{frame}{KL entre gaussianas (1/2) -- Partindo da definição}
    A KL entre duas distribuições quaisquer é:
    \vspace{0.3em}
    \[
      D_{KL}\!\left[q \,\|\, p\right]
      = \int q(\mathbf{z})\,\log\!\left(\frac{q(\mathbf{z})}{p(\mathbf{z})}\right) d\mathbf{z}
      = \mathbb{E}_{q}\!\left[\log q(\mathbf{z})\right]
      - \mathbb{E}_{q}\!\left[\log p(\mathbf{z})\right]
    \]
    \vspace{0.3em}
    \visible<2->{
      Para $q(\mathbf{z}|\mathbf{x},\bm{\phi}) = \mathcal{N}(\boldsymbol{\mu}, \boldsymbol{\Sigma})$ e
      $p(\mathbf{z}) = \mathcal{N}(0, \mathbf{I})$, o $\log$ de cada
      densidade gaussiana é uma forma quadrática:
      \[
        \log p(\mathbf{z}) = -\tfrac{D_z}{2}\log(2\pi) - \tfrac{1}{2}\,\mathbf{z}^{\top}\mathbf{z}
      \]
      \[
        \log q(\mathbf{z}) = -\tfrac{D_z}{2}\log(2\pi) - \tfrac{1}{2}\log|\boldsymbol{\Sigma}| - \tfrac{1}{2}(\mathbf{z}-\boldsymbol{\mu})^{\top}\boldsymbol{\Sigma}^{-1}(\mathbf{z}-\boldsymbol{\mu})
      \]
      \vspace{0.1em}
      \footnotesize Os termos $-\tfrac{D_z}{2}\log(2\pi)$ cancelam.
      No próximo slide usamos a identidade
      $\mathbb{E}_q[\mathbf{z}\mathbf{z}^{\top}] = \boldsymbol{\Sigma} + \boldsymbol{\mu}\boldsymbol{\mu}^{\top}$ (definição de
      covariância: $\boldsymbol{\Sigma} = \mathbb{E}_q[\mathbf{z}\mathbf{z}^{\top}] - \boldsymbol{\mu}\boldsymbol{\mu}^{\top}$).
    }
  \end{frame}

  %--- Slide: KL entre gaussianas 2/2 ---
  \begin{frame}{KL entre gaussianas (2/2) -- Forma fechada}
    Tomando esperanças sob $q$ (usando
    $\mathbb{E}_q[\mathbf{z}] = \boldsymbol{\mu}$ e $\mathbb{E}_q[\mathbf{z}\mathbf{z}^{\top}] = \boldsymbol{\Sigma} + \boldsymbol{\mu}\boldsymbol{\mu}^{\top}$):
    \vspace{0.3em}
    \begin{align*}
      \mathbb{E}_{q}\!\left[-\tfrac{1}{2}\mathbf{z}^{\top}\mathbf{z}\right]
      &= -\tfrac{1}{2}\!\left(\mathrm{Tr}[\boldsymbol{\Sigma}] + \boldsymbol{\mu}^{\top}\boldsymbol{\mu}\right) \\[0.2em]
      \mathbb{E}_{q}\!\left[-\tfrac{1}{2}(\mathbf{z}-\boldsymbol{\mu})^{\top}\boldsymbol{\Sigma}^{-1}(\mathbf{z}-\boldsymbol{\mu})\right]
      &= -\tfrac{1}{2}\,D_z
    \end{align*}
    \vspace{0.2em}
    Substituindo na definição da KL:
    \[
      \boxed{
        D_{KL}\!\left[\mathcal{N}(\boldsymbol{\mu},\boldsymbol{\Sigma})\,\|\,\mathcal{N}(0,\mathbf{I})\right]
        = \tfrac{1}{2}\!\left(
          \underbrace{\mathrm{Tr}[\boldsymbol{\Sigma}]}_{\text{dispersão}}
          + \underbrace{\boldsymbol{\mu}^{\!\top}\boldsymbol{\mu}}_{\text{desvio da origem}}
          - \underbrace{D_{z}}_{\text{norm}}
          - \underbrace{\log\det\boldsymbol{\Sigma}}_{\text{volume}}
        \right)
      }
    \]
    \note{
      Cada termo tem um papel pedagogicamente claro: $\mathrm{Tr}[\boldsymbol{\Sigma}]$
      penaliza variância grande, $\boldsymbol{\mu}^{\top}\boldsymbol{\mu}$ penaliza se afastar da
      origem, $\log\det\boldsymbol{\Sigma}$ penaliza variância pequena demais (evita
      colapso), $D_z$ é só normalização.
    }
  \end{frame}

  %--- Slide: Exemplo numérico KL ---
  \begin{frame}{Exemplo Numérico: Calculando o Termo KL}
    \footnotesize Imagine um VAE treinado em MNIST com latente 2D.
    Para a imagem de um ``7'' específico, o \emph{encoder} produz uma
    distribuição $q(\mathbf{z}\mid \mathbf{x}, \bm{\phi}) = \mathcal{N}(\boldsymbol{\mu}, \boldsymbol{\Sigma})$ sobre o
    latente, compacta o suficiente para ``codificar 7'', mas com alguma
    incerteza para suavizar o espaço. O termo KL mede quão longe essa
    distribuição está do prior $\mathcal{N}(0, \mathbf{I})$.

    \normalsize
    Encoder produz $\boldsymbol{\mu} = (0.5,\, -0.8)$ e
    $\boldsymbol{\Sigma} = \mathrm{diag}(0.25,\, 1.5)$ (latente 2D, $D_z = 2$).
    \begin{columns}[T]
      \begin{column}{0.50\textwidth}
        \textbf{Calculando cada termo:}
        \begin{align*}
          \mathrm{Tr}[\boldsymbol{\Sigma}] &= 0.25 + 1.5 = 1.75 \\
          \boldsymbol{\mu}^{\top}\boldsymbol{\mu} &= 0.25 + 0.64 = 0.89 \\
          \log\det\boldsymbol{\Sigma} &= \log(0.25) + \log(1.5) \\
                          &\approx -1.386 + 0.405 \\
                          &= -0.981 \\
          D_z &= 2
        \end{align*}
      \end{column}
      \begin{column}{0.48\textwidth}
        \begin{block}{Substituindo}
          \begin{align*}
            D_{KL} &= \tfrac{1}{2}\bigl(1.75 + 0.89 - 2 + 0.981\bigr) \\
                   &= \tfrac{1}{2}\,(1.621) \approx \mathbf{0.81}
          \end{align*}
        \end{block}
        Os dois termos que dominam:
        $\boldsymbol{\mu}^{\top}\boldsymbol{\mu} = 0.89$ (desvio da origem) e
        $-\log\det\boldsymbol{\Sigma} = +0.98$ (variância pequena $\sigma^{2}=0.25$
        é penalizada para evitar colapso).
      \end{column}
    \end{columns}
    \note{
      Note que o termo $-\log\det\boldsymbol{\Sigma}$ explode quando $\sigma_i^2 \to 0$:
      o KL impede que o encoder colapse a distribuição num único ponto.
      Esse é o efeito \emph{anti-colapso} do prior gaussiano isotrópico.
    }
  \end{frame}

  %--- Slide: Monte Carlo sample ---
  \begin{frame}{VAE -- Aproximação por \textit{Monte Carlo}}
    O primeiro termo do ELBO ainda é uma integral intratável. Porém,
    como é uma \textbf{esperança}, podemos aproximá-lo por
    \textit{Monte Carlo}:
    \vspace{0.3em}
    \[
      \int q(\mathbf{z} \mid \mathbf{x}, \bm{\phi})\,\log p(\mathbf{x} \mid \mathbf{z}, \thetav)\, d\mathbf{z}
      \approx \log p(\mathbf{x} \mid \mathbf{z}^{*}, \thetav),
      \quad \mathbf{z}^{*} \sim q(\mathbf{z} \mid \mathbf{x}, \bm{\phi})
    \]
    \vspace{0.3em}
    \begin{block}{ELBO aproximado}
      \[
        \mathrm{ELBO}[\thetav, \bm{\phi}] \approx
        \log p(\mathbf{x} \mid \mathbf{z}^{*}, \thetav)
        - D_{KL}\!\left[q(\mathbf{z} \mid \mathbf{x}, \bm{\phi})\,\|\,p(\mathbf{z})\right]
      \]
    \end{block}
    \note{
      Uma amostra só já dá um estimador não-enviesado. Na prática usa-se
      \emph{minibatch} + 1 amostra por instância, algo que o SGD tolera
      bem.
    }
  \end{frame}

  %======================================================================
  % SECTION 4: REPARAMETRIZATION TRICK
  %======================================================================
  \section{Reparametrization Trick}

  %--- Slide: reparametrization trick motivação ---
  \begin{frame}{\textit{Reparametrization Trick}}
    \begin{itemize}
      \item Como temos um \textbf{termo estocástico} no meio da rede
      neural (a amostragem $\mathbf{z}^{*} \sim q(\mathbf{z} \mid \mathbf{x}, \bm{\phi})$), é
      \emph{inviável} computar \textit{backpropagation} diretamente.
      \item \textbf{Solução}: reparametrizar a saída do \textit{encoder}
      de modo a tirar a amostragem do caminho do gradiente:
    \end{itemize}
    \vspace{0.3em}
    \[
      \mathbf{z}^{*} = \boldsymbol{\mu} + \boldsymbol{\Sigma}^{1/2}\,\boldsymbol{\epsilon}^{*},
      \qquad \boldsymbol{\epsilon}^{*} \sim \mathcal{N}(0, \mathbf{I})
    \]
    \vspace{0.2em}
    Para $\boldsymbol{\Sigma}$ diagonal (caso usual), $\boldsymbol{\Sigma}^{1/2}\,\boldsymbol{\epsilon}^{*} = \boldsymbol{\sigma} \odot \boldsymbol{\epsilon}^{*}$ (forma elemento a elemento).
    \begin{block}{Ideia chave}
      O ruído $\boldsymbol{\epsilon}^{*}$ é amostrado \emph{fora} da rede; $\boldsymbol{\mu}$ e
      $\boldsymbol{\Sigma}$ permanecem no caminho diferenciável.
    \end{block}
  \end{frame}

  %--- Slide: reparametrization diagrama ---
  \begin{frame}{\textit{Reparametrization Trick} -- Diagrama}
    \centering
    \resizebox{\linewidth}{!}{%
    \begin{tikzpicture}[
      every node/.style={font=\scriptsize},
      block/.style={draw, thick, minimum width=1.1cm, minimum height=0.9cm,
                    rounded corners=2pt},
      enc/.style={block, draw=EncRed, fill=EncRed!10, text=EncRed},
      dec/.style={block, draw=DecBlue, fill=DecBlue!10, text=DecBlue},
      reparam/.style={draw=NoiseAmber, dashed, thick, rounded corners=3pt,
                      fill=NoiseAmber!10, inner sep=5pt},
      arrow/.style={->, thick, >=stealth},
      ann/.style={font=\tiny, text=black!70, align=center},
      callout/.style={draw, dashed, rounded corners=1pt, fill=white,
                      inner sep=2pt, font=\tiny, align=center}
    ]
      % Main pipeline
      \node (x) {$\mathbf{x}$};
      \node[enc, right=0.4cm of x] (g) {$g[\mathbf{x},\bm{\phi}]$};
      \node[right=0.55cm of g] (params) {$\boldsymbol{\mu},\,\boldsymbol{\Sigma}^{1/2}$};
      \node[reparam, right=0.55cm of params]
        (rep) {$\mathbf{z}^{*} = \boldsymbol{\mu} + \boldsymbol{\Sigma}^{1/2}\boldsymbol{\epsilon}^{*}$};
      \node[dec, right=0.55cm of rep] (f) {$f[\mathbf{z}^{*},\thetav]$};
      \node[right=0.4cm of f] (px) {$p(\mathbf{x}\mid \mathbf{z}^{*},\thetav)$};

      % Noise source below
      \node[draw, NoiseAmber, fill=NoiseAmber!20, rounded corners=2pt,
            below=0.8cm of rep, minimum width=2.6cm, minimum height=0.7cm]
        (eps) {$\boldsymbol{\epsilon}^{*} \sim \mathcal{N}(0,\mathbf{I})$};

      \draw[arrow] (x) -- (g);
      \draw[arrow] (g) -- (params);
      \draw[arrow] (params) -- (rep);
      \draw[arrow, NoiseAmber] (eps) -- (rep);
      \draw[arrow] (rep) -- (f);
      \draw[arrow] (f) -- (px);

      % Labels under boxes
      \node[ann, below=0.08cm of x] {entrada};
      \node[ann, below=0.08cm of g] {encoder};
      \node[ann, below=0.08cm of params] {saída do encoder};
      \node[ann, below=0.08cm of f] {decoder};
      \node[ann, below=0.08cm of px] {verossimilhança};

      % Callouts above
      \node[callout, draw=NoiseAmber, text=NoiseAmber!80!black,
            above=0.9cm of rep]
        (note1) {ruído amostrado \emph{fora} da rede\\(determinístico em $\boldsymbol{\mu},\boldsymbol{\Sigma}$)};
      \node[callout, draw=EncRed, text=EncRed,
            above=0.9cm of g, xshift=-0.2cm]
        (note2) {gradiente flui até $\bm{\phi}$\\via $\boldsymbol{\mu}(\bm{\phi}),\,\boldsymbol{\Sigma}(\bm{\phi})$};
      \node[callout, draw=DecBlue, text=DecBlue,
            above=0.9cm of f, xshift=0.6cm]
        (note3) {gradiente flui até $\thetav$\\via $f(\mathbf{z}^{*},\thetav)$};

      \draw[->, dashed, NoiseAmber!80, thin] (note1) -- (rep);
      \draw[->, dashed, EncRed!70, thin] (note2) -- (g);
      \draw[->, dashed, DecBlue!70, thin] (note3) -- (f);
    \end{tikzpicture}}\\
    \vspace{0.3em}
    \tiny Adaptado de Prince (2023), Fig.~17.11.
  \end{frame}

  %======================================================================
  % SECTION 5: VAE EM AÇÃO
  %======================================================================
  \section{VAE em ação}

  %--- Slide: amostras e interpolação ---
  %  Figuras geradas por scripts/train_vae_mnist.py
  \begin{frame}{VAE em ação -- Amostras e interpolação}
    \begin{columns}[T]
      \begin{column}{0.48\textwidth}
        \centering
        \IfFileExists{slides/img/05-geradores/vae-samples.png}{%
          \includegraphics[width=0.72\linewidth]{slides/img/05-geradores/vae-samples.png}%
        }{%
          \fbox{\parbox[c][3.6cm][c]{0.85\linewidth}{\centering
            \textit{TODO: rode\\\texttt{scripts/train\_vae\_mnist.py}\\
            para gerar \texttt{vae-samples.png}.}}}%
        }\\
        \vspace{0.3em}
        \scriptsize \textbf{Amostragem}: $\mathbf{z}\sim p(\mathbf{z})$, depois
        $\mathbf{x} = f(\mathbf{z},\thetav)$. Resultado reconhecível mas \emph{borrado}
        (limite inferior não é apertado).
      \end{column}
      \begin{column}{0.48\textwidth}
        \centering
        \IfFileExists{slides/img/05-geradores/vae-interpolation.png}{%
          
\includegraphics[width=\linewidth]{slides/img/05-geradores/vae-interpolation.png}%
        }{%
          \fbox{\parbox[c][3.6cm][c]{0.85\linewidth}{\centering
            \textit{TODO: rode\\\texttt{scripts/train\_vae\_mnist.py}\\
            para gerar \texttt{vae-interpolation.png}.}}}%
        }\\
        \vspace{0.3em}
        \scriptsize \textbf{Interpolação}: transição \emph{suave} entre
        dois dígitos confirma que o latente é contínuo e semanticamente
        organizado.
      \end{column}
    \end{columns}
  \end{frame}

  %--- Slide: posterior collapse ---
  \begin{frame}{Patologia: \textit{posterior collapse}}
    \begin{columns}[T]
      \begin{column}{0.55\textwidth}
        \begin{itemize}
          \item Quando o decoder $f$ é \textbf{poderoso demais}, ele
          consegue modelar $\mathbf{x}$ \emph{ignorando} $\mathbf{z}$.
          \item O termo KL pressiona $q(\mathbf{z}\mid \mathbf{x},\bm{\phi})\to p(\mathbf{z})$:
          encoder colapsa para o prior. Latente vira \emph{inútil}.
          \item Sintomas durante o treino:
          \begin{itemize}
            \item $D_{KL}\to 0$;
            \item $\boldsymbol{\mu}(\mathbf{x})\to \mathbf{0}$, $\boldsymbol{\Sigma}(\mathbf{x})\to \mathbf{I}$ para todo $\mathbf{x}$.
          \end{itemize}
          \item Mitigações:
          \begin{itemize}
            \item \textbf{KL annealing}: rampa do peso do KL de $0$ a $1$;
            \item \textbf{Free bits}: troca $D_{KL}$ por
            $\sum_i \max(\lambda,\, D_{KL,i})$. Cada dimensão tem
            $\lambda$ nats ``grátis'' antes do KL pressionar;
            \item reduzir capacidade do decoder.
          \end{itemize}
        \end{itemize}
      \end{column}
      \begin{column}{0.43\textwidth}
        \centering
        \begin{tikzpicture}
          \begin{axis}[
            width=\linewidth, height=4.6cm,
            xlabel={\scriptsize época}, ylabel={\scriptsize valor},
            xtick=\empty, ytick=\empty,
            ymin=0, ymax=2.2, xmin=0, xmax=10,
            axis lines=left, no markers,
            legend style={font=\tiny,at={(0.98,0.98)},anchor=north east,draw=none,fill=none},
            domain=0:10, samples=80,
          ]
            \addplot[thick,EncRed] {1.5*exp(-0.6*x) + 0.05};
            \addlegendentry{KL}
            \addplot[thick,DecBlue] {1.8 - 0.6*exp(-0.4*x)};
            \addlegendentry{reconstrução}
          \end{axis}
        \end{tikzpicture}\\
        \vspace{0.2em}
        \scriptsize KL cai a zero enquanto a reconstrução estagna: o
        decoder ignora $\mathbf{z}$.
      \end{column}
    \end{columns}
  \end{frame}

  %======================================================================
  % PANORAMA
  %======================================================================
  \section{Panorama}

  %--- Slide: trilema ---
  \begin{frame}{Onde VAEs ficam no trilema?}
    \begin{columns}[T]
      \begin{column}{0.52\textwidth}
        \begin{itemize}
          \item \textbf{Amostragem rápida}: $\checkmark$ (1 forward do
          decoder).
          \item \textbf{Cobertura / diversidade}: $\checkmark$
          (maximização do ELBO cobre toda a distribuição).
          \item \textbf{Qualidade visual}: $\times$ (imagens
          tipicamente borradas; o limite inferior não é apertado).
        \end{itemize}
        \vspace{0.4em}
      \end{column}
      \begin{column}{0.44\textwidth}
        \centering
        \resizebox{\linewidth}{!}{\generativetrilemma[vaenf]}
      \end{column}
    \end{columns}
    \begin{block}{Posicionamento}
      VAEs ocupam o canto \emph{velocidade + cobertura} do trilema,
      ao lado dos \textit{Normalizing Flows}. Para qualidade,
      combine VAEs com modelos difusores, base da arquitetura
      \textit{Latent Diffusion}.
    \end{block}
  \end{frame}

  %======================================================================
  % RESUMO
  %======================================================================
  \section*{Resumo}

  \begin{frame}{Resumo}
    \begin{itemize}
      \item \textit{Autoencoders}: encoder $g$ e decoder $f$
      reconstruindo a entrada; latente útil para representação.
      \item \textit{Denoising AE}: entrada corrompida; rede aprende a
      remover ruído.
      \item Aplicação de AE: detecção de anomalias via erro de
      reconstrução.
      \item \textit{VAE}: \textit{autoencoder} probabilístico, latente
      é uma distribuição $q(\mathbf{z}\mid \mathbf{x},\bm{\phi})$.
      \item Identidade mestre:
      $\log p(\mathbf{x}) = \mathrm{ELBO} + D_{KL}[q\|p(\mathbf{z}\mid \mathbf{x})]$; bound
      apertado quando $q$ aproxima a posterior.
      \item ELBO $=$ reconstrução $-$ KL$[q(\mathbf{z}\mid \mathbf{x},\bm{\phi})\|p(\mathbf{z})]$
      (logo $-$ELBO $=$ reconstrução $+$ KL);
      para $p(\mathbf{x}\mid \mathbf{z})$ gaussiana, reconstrução $=$ MSE.
      \item \textit{Reparametrization trick} permite backprop através
      da amostragem.
      \item \textit{Posterior collapse}: KL $\to 0$, latente vira
      inútil; mitigado com KL annealing ou free bits.
    \end{itemize}
  \end{frame}

  %======================================================================
  % REFERÊNCIAS E LEITURAS
  %======================================================================
  \section*{Referências e Leituras}

  \begin{frame}{Leituras Recomendadas}
    \leiturasize
    \begin{itemize}
      \item Capítulo 17: \fullcite{prince2023understanding}
      \item \fullcite{kingma2014auto}
      \item \fullcite{vincent2008denoising}
      \item \fullcite{higgins2017beta}
    \end{itemize}
  \end{frame}

  \begin{frame}[allowframebreaks]{Referências}
    \sloppy
    \printbibliography[heading=none]
  \end{frame}

\end{document}
